\section{LLM 对话原理：上下文与``模拟对话''}
\label{sec:llm-dialogue}

\subsection{核心认知}

LLM（大语言模型）与生俱来的能力只有一项：\textbf{给定一段文本，预测下一个最合理的 token}。

所谓的``对话''、``记忆''、``上下文理解''全都是 prompt engineering 包装出来的产物。

\subsection{每次请求都发送完整历史}

每次用户发一条消息，Hermes（或其他 LLM 客户端）做的事：

\begin{lstlisting}
用户输入: "给我配飞书"

| Hermes 组装 prompt

[System Prompt (角色设定)]
[环境信息 (OS, CWD, ...)]
[持久化记忆 (user preferences)]
[历史对话 (压缩后)]
  - 用户: 配置飞书
  - 助手: 查日志 -> 发现 lark-oapi 缺失 -> 安装
  - 用户: 删掉错误 user
  - 助手: 改 .env -> 重启
[工具调用结果]
[本轮新消息: "给我配飞书"]

| 发送给 LLM API (几十万 token)

LLM 做文本接龙 -> 返回回复
\end{lstlisting}

用户只发了几个字节。客户端（Hermes）拼了几十万 token 发给 API。

\subsection{对话是模拟出来的}

\begin{longtable}{ll}
\toprule
\textbf{你以为的} & \textbf{实际发生的} \\
\midrule
\endhead
模型在``和你对话'' & 模型在做文本接龙 \\
模型``记得''你刚才说了什么 & 那段文本就在输入里 \\
对话是原生模式 & 对话是 prompt 格式的产物 \\
模型有``意图'' & 模型只是在模仿训练数据里的对话模式 \\
\bottomrule
\end{longtable}

关键创新是 \textbf{ChatML 格式}（由 OpenAI 提出）：

\begin{lstlisting}
user: 法国的首都是哪里？
assistant: 法国的首都是巴黎。
user: 那日本呢？
assistant: [LLM 在这里做文本接龙]
\end{lstlisting}

模型在训练时见过大量这种格式，学会了在这种格式下应该如何``扮演''。

\subsection{KV Cache 与 Prompt Caching}

每次发消息确实需要重新处理整个 prompt，但有两个优化：

\subsubsection{KV Cache（推理引擎层面）}

\begin{itemize}
  \item \textbf{Prefill 阶段}：把完整 prompt 一次性输入，计算所有 token 的 KV cache（主要计算瓶颈）
  \item \textbf{Decode 阶段}：每吐一个 token，只需要增量计算最后一个 token 的 KV cache，复用前面的
\end{itemize}

\subsubsection{Prompt Caching（API 层面）}

\begin{itemize}
  \item 无缓存：每次重新算全部历史的 KV，全部按 input token 计费
  \item 有缓存：重复部分按约 10\% 缓存费率，新内容按全价
\end{itemize}

DeepSeek、Claude、Gemini、OpenAI 等主流模型均支持。

\subsection{启示}

\begin{enumerate}
  \item 对话不是 LLM 的原生能力，而是通过数据格式和 prompt 设计模拟出来的
  \item 每次请求都发送完整历史（包括所有历史对话 + system prompt + 环境信息）
  \item 所谓``记忆力好''只是因为每次的输入里都包含了全部历史文本
  \item 用户端的流量很小（几字节），但到 API 的流量很大（几十万 token）
\end{enumerate}
